\begin{frame}{자주 하는 실수: 정의를 건너뛰고 알고리즘부터}
  \begin{itemize}
    \item 사례: 스타트업이 ``이탈 예측 모델 만들자''는 목표만 갖고
      ``\textbf{그럼 신경망을 쓸까, XGBoost를 쓸까?}''부터 논의.
    \item 몇 주 뒤 모델은 완성됐지만 \textbf{성능이 엉망} ---
      ``이탈''의 정의(30일 미접속? 구독 해지? 환불 요청?)조차
      \textbf{팀마다 다르게} 이해하고, $y$(라벨) 자체가 일관성 없이
      매겨져 있었음.
  \end{itemize}
  \vskip0.5em
  \kq{알고리즘 선택에 쓴 시간보다, $y$ 를 정의하는 데 쓴 시간이 훨씬
  적었던 것이 진짜 원인.}
  \vskip0.5em
  \begin{itemize}
    \item ``어떤 모델이 가장 정확할까''는 \textbf{프로젝트 중반 이후}의
      질문. 초반에 가장 먼저 답해야 할 질문은 ``$x$ 는 무엇이고, $y$ 는
      정확히 무엇이며, 그 정의에 \textbf{모두가 동의하는가}''.
    \item Week 8, 16 팀 프로젝트: 데이터를 내려받자마자 모델링에
      뛰어들지 말고, 이 세 가지 질문에 \textbf{팀 전체가 명시적으로}
      답을 적어두기.
  \end{itemize}
\end{frame}
